% ============================================================
%  H. Brøchner, "Problemet om Tro og Viden" (1868)
%  Preface (pp. 1–7)
%  TRANSLATION (English)
% ============================================================
\documentclass[12pt,a4paper]{article}
\usepackage[utf8]{inputenc}
\usepackage[T1]{fontenc}
\usepackage[english]{babel}
\usepackage{microtype}
\usepackage{setspace}
\usepackage[margin=1.2in]{geometry}
\usepackage{fancyhdr}
\usepackage{hyperref}

\hypersetup{
  pdftitle={The Problem of Faith and Knowledge — Preface},
  pdfauthor={H. Brøchner},
  hidelinks
}

\pagestyle{fancy}
\fancyhf{}
\rhead{Brøchner, \emph{The Problem of Faith and Knowledge} (1868)}
\cfoot{\thepage}

\setstretch{1.3}

\title{\textbf{The Problem of Faith and Knowledge}\\[6pt]
  \large A Historical-Critical Treatise\\[10pt]
  \normalsize Preface}
\author{H.\ Brøchner\\[4pt]
  \small Translated from the Danish by Hans Halvorson}
\date{Copenhagen: P.\ G.\ Philipsen, 1868}

\begin{document}
\maketitle
\thispagestyle{fancy}

\bigskip

\noindent This work reproduces, in all essential respects unchanged,
the content of the lectures I delivered at the University of Copenhagen
in the autumn semester of 1867 on the relation between faith and
knowledge, between religion, philosophy, and ethics. These lectures
encompassed the historical-critical part of the investigation and
formed the preparation for an attempt at an independent resolution of
the problem under discussion, which will be published in the course of
this year.

The historical-critical part of the investigation provides, beyond a
brief sketch of the problem's genesis, a detailed critique of those
proposed solutions that have acquired a special significance in our
literature --- among them one, the solution originating with
S.\ Kierkegaard, that was executed with an originality, a seriousness,
and a depth that will preserve for it this significance at all times.
Through the critique of these proposed solutions --- which, in the
course of their development and justification, have stood in a
sustained polemical or mediating relation to earlier and foreign
contributions, and which can partly be regarded as results of
previously dominant tendencies, and have partly substantially
accomplished the critique of such tendencies --- the most significant
of the older and contemporary proposals have been illuminated, and in
part the complete preconditions have been supplied for a final judgment
upon them. This historical critique will then find its supplement in
the work that attaches itself to this one.

The present state of the problem in our literature is marked by the
conflict between, on the one side, a form of theology that borrows in
an eclectic fashion from newer philosophical hybrid formations the
weapons with which it will defend the union of faith and knowledge,
and, on the other side, the anti-theological-religious standpoint,
which, by means of the principled opposition between faith and
knowledge, aims to secure the validity of religion and science in
their mutual independence and the possibility of their peaceful
coexistence within consciousness --- while at the same time, by virtue
of this same principled opposition, denying theology all right to bear
the name of science, a name theology claims for itself by demanding to
be a ``science of faith'' [\textit{Troesvidenskab}].

The philosophizing theology, which maintains the possibility of
appropriating and expressing the content of faith in the form of
knowledge, and which therefore not only wishes to claim the name of
science for theology but assigns to it, as the concluding form of
revelatory science, primacy among the sciences and judicial authority
over them, has had as its historical point of departure the speculative
system that, through an ambiguous extension of the concept of
``religion,'' brought religion in its specific sense into an identical
relation with philosophy, from which it was to be distinguished not by
any difference of content, but only by a difference of form. After
this ambiguous unity had to be abandoned in the struggles that
developed following the emergence of Hegel's philosophy of religion,
the theology influenced by philosophy, exchanging the designation of
speculation for that of the Christian, nevertheless preserved the
tendency toward mediation; but it now rests primarily upon a different
philosophical orientation --- and in particular upon the system that,
by introducing a redoubling into philosophy, creates room for a ``free
thinking'' which borrows from its close relative, the imagination, the
``magic cloak'' [\textit{Fjederham}] enabling it to soar above the
abysses of mystery; and it cautiously inserts a qualification that cuts
off even the ``higher science's'' consequences whenever thought,
despite its ``freedom,'' involuntarily comes to follow its own
essential laws.

The anti-theological theory, which likewise presupposes Hegelian
philosophy, but whose relation to that presupposition is originally
negative and polemical, had as its historical incitement also the
protest --- lodged by a philosophy that presents itself at once as the
consequence of and as the opposition to that speculative doctrine of
mediation --- against the unity of religion and philosophy. The most
significant representative of this tendency is Feuerbach --- a man
whose writings are known only very incompletely and judged only very
superficially by those in our literature who have criticized him.
Conceiving religion's standpoint as practical in opposition to
science's theoretical standpoint, and determining the former ever more
sharply through the emphasis on its alogical element, he arrived at
the irreconcilable opposition as his view of the relation, and
demanded, in the interest of genuine cognition and genuine ethical
action, that religion be absorbed into passion and ethics. Religion as
such was to disappear --- as something alogical, as an imperfect
transitional stage in the development of the human spirit's
self-consciousness --- leaving behind only its name in an improper and
figurative sense.

In agreement with Feuerbach regarding the opposition, but disagreeing
as to its grounds and its application, S.\ Kierkegaard draws the
distinction between religion and philosophy in religion's interest.
The human spirit, which in the restlessness of existence cannot grasp
the Eternal in the objective certainty of cognition, and which
acknowledges the limits of its own cognition, finds only in the
religious relation, in faith, a subjective existential certainty
regarding what is its deepest interest; it finds there, and only
there, its true reconciliation. The validity of knowledge is not
denied, but objective knowledge is bounded to a domain in which the
spirit's deepest interests cannot find their satisfaction. Faith and
knowledge thus each receive their own territory, and through this, and
through knowledge's subordination, the compatibility of both within
the same consciousness is to become possible.

This Kierkegaardian conception of the problem of the relation between
faith and knowledge has been taken up by Professor R.\ Nielsen and
presented in such a way that some have wished to place him on a par in
originality with the theory's originator, attributing to him the
resolution of a task of the same creativity as the one Kierkegaard
strove to resolve. How little this view is warranted will be
demonstrated in this work. It will emerge that the problem has not
been advanced by Professor Nielsen; that the form he has given the
solution has introduced only uncertainty and vacillation into the
decisive determinations; that the attempt at a psychological and
metaphysical grounding of the separation of the principle of faith
from the principle of knowledge as absolutely heterogeneous with it
has failed and leads to irresolvable contradictions; and that the same
holds for his grounding of the proposition that the absolutely
heterogeneous principles can be united within the same consciousness.

When, then, alongside S.\ Kierkegaard --- the genuine representative
of the anti-theological-religious standpoint --- such detailed
attention is devoted to Professor Nielsen, it is not because the
theory has received from him a deeper grounding or a richer
development, but because it has through him exchanged the character of
an esoteric doctrine for an exoteric one, penetrated into the broad
circles of ``the cultivated,'' and in these circles is in the process
of changing from a question of knowledge into a question of faith, and
of being fixed as dogma by virtue of an unappealable
\textit{ipse dixit}.\footnote{Brøchner's text has the Greek:
$\alpha\dot{\upsilon}\tau\grave{o}\varsigma\ \check{\varepsilon}\varphi\eta$
--- the Pythagorean formula, ``He himself said it,'' used to denote
unquestioned reliance on authority.}

As representative of the philosophizing theology, Bishop Martensen
appears in this work; but since the argument is closely tied to what
is given in his writings, the universal, the typical in the standpoint
is nowhere lost from view, nor forgotten in favor of the individual.
It has not been necessary to give direct consideration to other
theological attempts in our literature.

The ultimate task, whose resolution is prepared through the critique
--- which immediately can yield only a negative result --- is the
vindication of a standpoint that has not found expression in our
discussions, and that has been advanced elsewhere only in forms that
cannot satisfy. This standpoint maintains such an opposition between
science --- philosophy in particular --- and the positive forms of
religion, between knowledge and faith understood as revelatory faith,
that knowledge and this faith cannot be brought to unity within the
same consciousness; --- it maintains the impossibility of a genuine
and substantial ethics, one that encompasses the whole, real human
life, developing on the basis of positive religion with its
supernatural law of action; --- and it maintains the inner,
indissoluble unity between genuine cognition and genuine ethical
action, the necessary dependence of the latter upon the former. But
while it maintains this opposition, it asserts with equal
decisiveness, on the other side --- distinguishing the essential
determination of the religious from the phenomenal forms of the
positive religions, in which it sees only a partial, imperfect,
finitizing expression of the former, and thus distinguishing what is
religious in the religions from the religions in their historical
forms --- the inner unity of the religious, the philosophical, and the
ethical, and through this inner unity the inner unity of human life:
the condition for its true and complete reconciliation, the
reconciliation within reality, the reconciliation of undivided human
nature with itself in its essential ground. In affirming the
essentiality and necessity of the religious relation for the human
being, whose self-consciousness is first fully actualized therein,
through God-consciousness, and who according to the fundamental
relation of its nature has never been and never will be without
religion; --- in further seeing ethical action and genuine cognition as
gaining a higher significance through their unity with the religious
fundamental relation; --- and in allowing genuine knowledge and genuine
ethical action to become an adequate form of actuality for the
religious, it conceives this knowledge and this action as containing
within themselves a Highest and Unconditioned, in which what is
essential in the religions is preserved. This standpoint we designate
as \textit{the religiously reconciled humane consciousness}.

It differs from the position occupied by Feuerbach not merely in that
its metaphysical fundamental concept is a different one --- one upon
which a cognition and an action bearing the stamp of ideality and
universal validity can be built, and from which the concept of human
freedom can be vindicated and a religious relation grounded to the
Absolute, with which the human being is in essential unity --- but
also in that it combats theology and dogma not only in the interest of
ethics and science, but first and foremost in the interest of religion
itself, in the interest of the infinity of the religious relation.

It differs from the old Hegelian school's conception of religion's
relation to cognition, in that it does not see the religious merely as
a form of knowledge, but as the expression of a fundamental relation
encompassing the whole of human nature in its concreteness, and in
that it does not let the medium of spiritual life be an illusory
abstract element of eternity, but reality with its real conditions and
limits.

With the anti-theological-religious theory it concurs in its
conception of theology's relation to science. Since it rests upon
scientific cognition and, in undivided human nature, assigns primacy
to cognition, it must be its task to vindicate the clear concept of
science and science's autonomy, grounded in its own nature, against
equivocations and against attempts, based on equivocations, to
subordinate science to the dominion of a foreign, heterogeneous
authority. But while there is thus, in an essential point, an
agreement between it and that theory, it enters into a decided
opposition to the latter when the relation between faith and
knowledge, between religion, philosophy, and ethics, is to be
determined. This opposition concerns Kierkegaard's conception, with
respect to which it will be shown that it undermines the certainty of
all cognition and allows the ethical to disappear into an abstract and
reality-less religious sphere. But it concerns even more precisely the
form in which the theory has been elaborated and presented by Nielsen
--- a form that Kierkegaard would certainly not recognize as his own.
The form in which Nielsen develops the theory leads, despite all the
more or less adroit turns that are made in order to avoid this
consequence, to an actual dualism in human nature, to a division of
the life of the spirit from its very ground by means of the ``absolutely
heterogeneous'' principles, to the inevitable self-dissolution of
consciousness; it corrupts the ethical at its root and opens the path
to the religious for all that is superstitious and absurd. With the
conviction that the absolute heterogeneity Nielsen posits between the
principle of knowledge on the one hand and the principles of faith and
action on the other is ``a misunderstanding that is ruinous for
religiosity, morality, and all genuine character development,'' I have
subjected his theory to a thoroughgoing critique, and sought to
demonstrate its untenability, its deficient grounding, its
self-contradictions, its ruinous consequences.

In publishing this contribution to the investigation of a problem that
through the centuries has had, and will continue to have, the most
decisive significance for the human spirit, it is my hope that it may
occasion a renewed, serious, and conscientious discussion --- a sharp
and open conflict between principles. Such a conflict, even if carried
on with sharp weapons, can only serve the interest of the cause; and
one must therefore hold dear the combatant for whom the cause, and not
some egoistic self-regard, is the highest consideration. A firm and
deep personal conviction I shall always respect; by sound reasons I
shall always be just as open to influence as I intend to be impervious
to anything other than these.

\end{document}
